\exercisesection{8}

\begin{enumerate}
\def\labelenumi{\arabic{enumi}.}
\item
  Let K be a field and A = K{[}{[}X, Y, Z{]}{]} the ring of formal power series in three indeterminates over K; let \(v'\) (resp. \(v''\) ) be the valuation on A with values in the group \(Z \times Z\) ordered lexicographically, such that \(v'(X) = (1, 0)\) , \(v'(Y) = (0, 1)\) , \(v'\) is improper in K{[}{[}Z{]}{]} (resp. \(v''(X) = (1, 0)\) , \(v''\) is improper in K{[}{[}Y{]}{]}, \(v''(Z) = (0, 1)\) ). Let \(\sigma\) be the automorphism of A leaving invariant K and X and such that \(\sigma(Y) = Z\) , \(\sigma(Z) = Y\) ; if B is the subring of A consisting of the elements invariant under \(\sigma\) and E (resp. F) the field of fractions of A (resp. B), the valuations \(v'\) and \(v''\) (extended canonically to E) have the same restriction v to F, F is complete with the topology defined by v and E is a quadratic extension of F; the two valuations \(v'\) , \(v''\) on E are dependent but not equivalent.
\item
  Let \(K_{0}\) be the field obtained by adjoining to the 2-adic field \(Q_{2}\) the roots of all the polynomials \(X^{2^{n}} - 2\) ; let v be the unique valuation on \(K_{0}\) extending the 2-adic valuation and let K be the completion of \(K_{0}\) with respect to v. Show that the polynomial \(X^{2} - 3\) is irreducible in K{[}X{]}; let \(K'\) be the field of roots of this polynomial and let \(v'\) be the extension of v to \(K'\) ; then
\end{enumerate}

\[
n = \left[ \mathrm{K} ^ {\prime}: \mathrm{K} \right] = 2, \quad e (v ^ {\prime} / v) = f (v ^ {\prime} / v) = 1.
\]

\begin{enumerate}
\def\labelenumi{\arabic{enumi}.}
\setcounter{enumi}{2}
\tightlist
\item
  Let k be the field of rational functions \(\mathbf{F}_{p}(\mathbf{X}_{n})_{n\in\mathbf{N}}\) in an infinity of indeterminates over the prime field \(F_{p}\) and let \(K = k(\mathbf{U}, \mathbf{V})\) be the field of rational functions in two indeterminates over k.
\end{enumerate}

\begin{enumerate}
\def\labelenumi{(\alph{enumi})}
\item
  Show that the element \(\mathbf{P}(\mathbf{U}) = \sum_{n=0}^{\infty} \mathbf{X}_n^p \mathbf{U}^{np}\) of the field of formal power series \(k((\mathbf{U}))\) is not algebraic over the field \(k(\mathbf{U})\) of rational functions. The mapping \(\mathbf{F}(\mathbf{U}, \mathbf{V}) \mapsto \mathbf{F}(\mathbf{U}, \mathbf{P}(\mathbf{U}))\) of \(k[\mathbf{U}, \mathbf{V}]\) to \(k((\mathbf{U}))\) can be extended to an isomorphism of \(\mathbf{K}\) onto a subfield of \(k((\mathbf{U}))\) ; the restriction to this subfield of the valuation on \(k((\mathbf{U}))\) equal to the order of the formal power series (§ 3, no. 3, Example 3) is a discrete valuation \(v\) on \(\mathbf{K}\) , whose residue field is \(k\) .
\item
  Let \(\mathbf{K}'\) be the algebraic extension \(\mathbf{K}(\mathbf{V}^{1/p})\) of \(\mathbf{K}\) so that \([\mathbf{K}':\mathbf{K}] = p\) ; if \(v'\) is the unique valuation on \(\mathbf{K}'\) extending \(v\) , show that \(e(v'/v) = f(v'/v) = 1\) . The ring of the valuation \(v'\) is therefore not a finitely generated module over the ring of the valuation \(v\) .
\end{enumerate}

\begin{enumerate}
\def\labelenumi{\arabic{enumi}.}
\setcounter{enumi}{3}
\item
  Let \(k\) be a field, \(\mathbf{K} = k(\mathbf{X}, \mathbf{Y})\) the field of rational functions in two indeterminates over \(k\) and \(v\) the valuation on \(\mathbf{K}\) with values in the group \(\mathbf{Z} \times \mathbf{Z}\) ordered lexicographically, such that \(v(\mathbf{X}) = (0, 1)\) , \(v(\mathbf{Y}) = (1, 0)\) . Let \(\mathbf{K}'\) be the field \(\mathbf{K}(\sqrt{\mathbf{X}})\) ; show that \(v\) has a single extension \(v'\) to \(\mathbf{K}'\) and that \(e(v'/v) = 2\) \(f(v^{\prime}/v)=1\) , but the ring of the valuation \(v^{\prime}\) is not a finitely generated module over the ring of the valuation \(v\) .
\item
  Let K be a field, A a valuation ring of K and L a finite algebraic extension of K. Let \(A' \supset A\) be another valuation ring of K; let \(A_i' (1 \leqslant i \leqslant m)\) be the valuation rings of L such that \(A_i' \cap K = A'\) and let \(k_i'\) be their respective residue fields. Let k be the residue field of \(A'\) and \(A''\) the valuation ring of K the canonical image of A; finally let \(A_{ij}' (1 \leqslant j \leqslant n_i)\) be the valuation rings of \(k_i'\) such that \(A_{ij}'' \cap k = A''\) and \(A_{ij}\) the inverse images in \(A_i'\) of the \(A_{ij}'' (1 \leqslant i \leqslant m, 1 \leqslant j \leqslant n_i)\) .
\end{enumerate}

\begin{enumerate}
\def\labelenumi{(\alph{enumi})}
\item
  Show that the \(\mathbf{A}_{ij}\) are valuation rings of \(\mathbf{L}\) which are distinct and such that \(\mathbf{A}_{ij} \cap \mathbf{K} = \mathbf{A}\) and every valuation ring \(\mathbf{B}\) of \(\mathbf{L}\) such that \(\mathbf{B} \cap \mathbf{K} = \mathbf{A}\) is equal to one of the \(\mathbf{A}_{ij}\) .
\item
  Show that
\end{enumerate}

\[
e \left(\mathrm{A} _ {i j} / \mathrm{A}\right) = e \left(\mathrm{A} _ {i} ^ {\prime} / \mathrm{A} ^ {\prime}\right) e \left(\mathrm{A} _ {i j} ^ {\prime \prime} / \mathrm{A} ^ {\prime \prime}\right) \quad \text { and } \quad f \left(\mathrm{A} _ {i j} / \mathrm{A}\right) = f \left(\mathrm{A} _ {i j} ^ {\prime \prime} / \mathrm{A} ^ {\prime \prime}\right).
\]

¶ 6. Let K be a field, v a valuation on K and A the ring of v.

\begin{enumerate}
\def\labelenumi{(\alph{enumi})}
\item
  Suppose that A is Henselian (Chapter III, § 4, Exercise 3) in which case we also say, by an abuse of language, that K is Henselian for v. Then, for every algebraic extension L of K and every valuation \(v'\) on L extending v, the ring \(A'\) of \(v'\) is Henselian (Chapter III, § 4, Exercise 4).
\item
  If K is complete with v and v is of height 1, A is Henselian. Give an example where K is complete with v and v is of height 2, but A is not Henselian (cf.~Exercise 1).
\item
  If K is linearly compact with v, show that K is Henselian. (In the notation condition (H) of Chapter III, § 4, Exercise 3, let L denote the set of prime ideals of A (totally) ordered by inclusion and \(\mathcal{L}\) the set of ordered pairs (B, φ), where B is a well ordered subset (with respect to the relation ⊃) of L and φ: p ↦ (Q\_p, Q'\_p) a mapping from B to the set A{[}X{]} × A{[}X{]}, with the following properties: (1) Q\_p and Q'\_p are monic of respective degrees deg( \(\overline{Q}\) ) and deg( \(\overline{Q}'\) ); (2) if p ⊃ q are two elements of B, the coefficients of the polynomials Q\_p - Q\_q and Q'\_p - Q'\_q belong to p; (3) the coefficients of P - Q\_pQ'\_p belong to p; (4) \(\bar{f}(Q_{p}) = \overline{Q}, \bar{f}(Q'_{p}) = \overline{Q}'\) . Define on L an ordering by setting (B, φ) ≤ (B', φ') when B ⊂ B' and φ' is an extension of φ to B'. Prove that L admits a maximal element (B₀, φ₀) and that B₀ has a last element equal to (0): obtain a contradiction from the opposite by considering two cases, according to whether or not B₀ has a last element; in the former case, if p is this last element, consider an element c ∈ A such that v(c) is the least value taken by v on the set of coefficients of P - Q\_pQ'\_p and the least prime ideal p' of A containing c; if p' = p, use Exercise 2 (a) of § 4 and the fact that a quotient of a linearly compact module by a closed submodule is linearly compact and hence complete. If on the contrary B₀ has no last element, use directly the hypothesis that A is linearly compact.)
\item
  In the notation of Exercise 5, show that for A to be Henselian, it is necessary and sufficient that \(\mathbf{A}'\) and \(\mathbf{A}''\) be so. (Note that, if \(\mathbf{P} \in \mathbf{A}'[\mathbf{X}]\) , there exists \(a \in \mathbf{A}\) such that \(a^n\mathbf{P}(\mathbf{X}/a) \in \mathbf{A}[\mathbf{X}]\) (where \(n\) is the degree of \(\mathbf{P}\) ); to verify axiom (H) of Chapter III, § 4, Exercise 3 for \(\mathbf{A}'\) , attention may be confined to the polynomials of \(\mathbf{A}[\mathbf{X}]\) ; use the fact that \(\mathbf{A}' = \mathbf{A}_{\mathfrak{p}}\) , where \(\mathfrak{p}\) is a prime ideal of \(\mathbf{A}\) contained in the maximal ideal \(\mathfrak{m}\) , and note that, if two polynomials in \((\mathbf{A}/\mathfrak{p})[\mathbf{X}]\) are strongly relatively prime, so are their images in \((\mathbf{A}/\mathfrak{m})[\mathbf{X}]\) .)
\item
  Suppose that A is Henselian; then, for every algebraic extension L of K, there exists (up to equivalence) only one valuation on L extending v. (If \(P \in A[X]\) is a monic irreducible polynomial, \(x_{i}\) ( \(1 \leqslant i \leqslant m\) ) the distinct roots of P in its splitting field N, show that, for every valuation \(v'\) on N extending v, the \(v'(x_{i})\) are all equal for \(1 \leqslant i \leqslant m\) .) Conversely, if the valuation ring A has this property, it is Henselian (observe that, in the integral closure \(A'\) of A in L, there can only exist one maximal ideal lying above m (Chapter V, § 1, Exercise 13) and conclude that, if \(P \in A[X]\) is a monic irreducible polynomial its image in \((\mathbf{A}/\mathfrak{m})[\mathbf{X}]\) cannot be a product of two relatively prime polynomials).
\end{enumerate}

\begin{enumerate}
\def\labelenumi{\arabic{enumi}.}
\setcounter{enumi}{6}
\tightlist
\item
  Let K be a field, v a valuation on K, A the ring of v and m the ideal of v. Suppose that A is Henselian; let L be a finite extension of K of degree n; let \(v'\) be the unique valuation on L extending v (Exercise 6 (e)), A' its ring and \(m'\) its ideal. Let x be any element of A'; show that the degree over A/m of the class \(\bar{x}\) of x in A'/m' is a divisor of n and that the order of the class of \(v'(x)\) in \(\Gamma_{v'}/\Gamma_v\) is a divisor of n (consider the minimal polynomial of x over K).
\end{enumerate}

¶ 8. Let K' be a field, B the ring of a valuation v on K', G a finite group of automorphisms of K', K the subfield of K' consisting of the elements invariant under G and A = K ∩ B, which is a valuation ring of K, corresponding to the valuation w = v\textbar K.

\begin{enumerate}
\def\labelenumi{(\alph{enumi})}
\item
  The valuations on \(\mathbf{K}'\) which extend \(w\) are the \(v \circ \sigma\) , where \(\sigma \in \mathcal{G}\) (no. 6, Corollary 1 to Proposition 6) and the integral closure \(\mathbf{A}'\) of \(\mathbf{A}\) in \(\mathbf{K}'\) is the intersection of the \(\sigma \cdot \mathbf{B}\) , where \(\sigma\) runs through \(\mathcal{G}\) . If \(\mathfrak{p}(\mathbf{B})\) is the intersection of \(\mathbf{A}'\) and the maximal ideal \(\mathfrak{m}(\mathbf{B})\) , then \(\sigma \cdot \mathfrak{p}(\mathbf{B}) = \mathfrak{p}(\sigma \cdot \mathbf{B})\) . Show that the decomposition group \(\mathcal{G}^{\mathbb{Z}}(\mathfrak{p}(\mathbf{B}))\) (Chapter V, § 2, no. 2) is the subgroup of \(\mathcal{G}\) consisting of the \(\sigma\) such that \(\sigma \cdot \mathbf{B} = \mathbf{B}\) ; write \(\mathcal{G}^{\mathbb{Z}} = \mathcal{G}^{\mathbb{Z}}(\mathfrak{p}(\mathbf{B}))\) and denote by \(\mathbf{K}^{\mathbb{Z}}\) the subfield of \(\mathbf{K}'\) consisting of the elements invariant under \(\mathcal{G}^{\mathbb{Z}}\) and by \(\mathbf{B}^{\mathbb{Z}}\) the ring of the valuation induced on \(\mathbf{K}^{\mathbb{Z}}\) by \(v\) , which is equal to \(\mathbf{B} \cap \mathbf{K}^{\mathbb{Z}}\) ; recall that the residue fields of \(\mathbf{B}^{\mathbb{Z}}\) and \(\mathbf{A}\) are the same and that the maximal ideal \(\mathfrak{m}(\mathbf{B}^{\mathbb{Z}})\) is generated by \(\mathfrak{m}(\mathbf{A})\mathbf{B}^{\mathbb{Z}}\) (Chapter V, § 2, no. 2, Proposition 4).
\item
  Let \(v^{\mathbf{Z}}\) denote the restriction of \(v\) to \(\mathbf{K}^{\mathbf{Z}}\) and \(\Gamma\) and \(\Gamma^{\mathbf{Z}}\) the order groups of \(w\) and \(v^{\mathbf{Z}}\) . Show that \(\Gamma = \Gamma^{\mathbf{Z}}\) . (Let \(\mathcal{G}^{\mathrm{D}} \supset \mathcal{G}^{\mathbf{Z}}\) be the subgroup of \(\mathcal{G}\) consisting of the \(\sigma \in \mathcal{G}\) such that \(\sigma\) . B is dependent on B (§ 7, no. 2). Argue by induction on the order of \(\mathcal{G}\) . If \((\mathcal{G}: \mathcal{G}^{\mathrm{D}}) > 1\) , let \(\mathbf{K}''\) be the field of invariants of \(\mathcal{G}^{\mathrm{D}}\) and \(\mathbf{A}'' = \mathbf{K}'' \cap \mathbf{B}\) ; show that the order group of \(v|\mathbf{K}''\) is equal to \(\Gamma\) : for this, use the Approximation Theorem (§ 7, no. 2, Theorem 1) in order to prove that, for all \(x \in K''\) , there exists \(y \in K''\) such that \(v(y) = v(x)\) and \(v(y_i) = 0\) for the conjugates \(y_i\) of \(y\) with respect to \(K\) , distinct from \(y\) ; then \(G\) may be replaced by \(G^D\) and the induction hypothesis may be applied. If \(G^D = G\) , let \(B_0\) be the valuation ring of \(K'\) generated by the \(\sigma.B\) , where \(\sigma \in G\) : let \(\bar{K}' = \kappa(B_0)\) be its residue field, \(\pi: B_0 \to \bar{K}'\) the canonical homomorphism, \(\bar{B} = \pi(B)\) a valuation ring of \(\bar{K}'\) and \(\bar{v}\) the corresponding valuation on \(\bar{K}'\) such that \(v = \bar{v} \circ \pi\) in \(B_0\) , so that the order group \(\bar{\Delta}\) of \(\bar{v}\) is a subgroup of the order group \(\Delta\) of \(v\) ; if \(\Delta_0 = \Delta/\bar{\Delta}\) and \(\bar{\omega}: \Delta \to \Delta_0\) is the canonical homomorphism, \(v_0 = \bar{\omega} \circ v\) is a valuation on \(K'\) corresponding to \(B_0\) and invariant under \(G\) . By taking quotients, \(G\) defines a group of automorphisms \(\bar{G}\) of \(\bar{K}'\) ; let \(N\) be the kernel of the canonical homomorphism \(G \to \bar{G}\) ; show that \(N \subset G^Z\) and deduce that, if \(N \neq \{e\}\) , \(G\) may be replaced by \(G/N\) and the induction hypothesis may be applied. Suppose finally that \(N = \{e\}\) ; let \(A_0 = K \cap B_0\) and \(w_0\) be the restriction of \(v_0\) to \(K\) ; show that the order group of \(w_0\) is \(\Delta_0\) and that \(\pi(A_0) = \bar{K}\) is the field of invariants of \(\bar{G}\) (cf.~no. 1, Lemma 2); if \(\bar{w}\) is the restriction to \(\bar{K}\) of \(\bar{v}\) and \(\bar{\Gamma}\) its order group, \(\Delta_0\) is canonically isomorphic to \(\Gamma/\bar{\Gamma}\) . On the other hand, the group \(G^Z(p(\bar{B}))\) is the canonical image of \(G^Z(p(B))\) in \(G^Z\) ; observe that \(G^D \neq G^Z\) and that the induction hypothesis may therefore be applied to prove that \(\bar{\Gamma}\) is equal to the order group \(F^Z(θ)\) of the restriction of \(\bar{v}\) to \(K^Z\). )
\item
  Let \(\mathcal{G}^{\mathrm{T}} = \mathcal{G}^{\mathrm{T}}(\mathfrak{p}(\mathrm{B}))\) be the inertia group of \(\mathfrak{p}(\mathrm{B})\) (Chapter V, § 2, no. 2) and let \(\mathbf{K}^{\mathrm{T}}\) denote the subfield of \(\mathbf{K}'\) consisting of the invariants of \(\mathcal{G}^{\mathrm{T}}\) and \(\mathbf{B}^{\mathrm{T}}\) the ring of the valuation induced on \(\mathbf{K}^{\mathrm{T}}\) by \(v\) , which is equal to \(\mathbf{B} \cap \mathbf{K}^{\mathrm{T}}\) ; recall that, if \(k\) and \(k'\) are the residue fields of A and B, the residue field \(k^{\mathrm{T}}\) of \(\mathbf{B}^{\mathrm{T}}\) is the greatest separable extension of \(k\) contained in the quasi-Galois extension \(k'\) of \(k\) and that \(\mathcal{G}^{\mathrm{Z}} / \mathcal{G}^{\mathrm{T}}\) is canonically isomorphic to the group of \(k\) -automorphisms of \(k'\) (or \(k^{\mathrm{T}}\) ) (Chapter V, § 2, Theorem 2 and Proposition 5). Let \(v^{\mathrm{T}}\) be the restriction of \(v\) to \(\mathbf{K}^{\mathrm{T}}\) ; show that the order group of \(v^{\mathrm{T}}\) is also equal to \(\Gamma\) (apply Lemma 2 of no. 1 to the extension \(\mathbf{K}^{\mathrm{T}}\) of \(\mathbf{K}^{\mathrm{Z}}\) ).
\end{enumerate}

\begin{enumerate}
\def\labelenumi{\arabic{enumi}.}
\setcounter{enumi}{8}
\tightlist
\item
  Let K be a field, L a finite algebraic extension of K, \(v'\) a valuation on L, A' its ring, v the restriction of \(v'\) to K and \(A = A' \cap K\) its ring.
\end{enumerate}

\begin{enumerate}
\def\labelenumi{(\alph{enumi})}
\item
  Suppose that A is Henselian (Exercise 6); show that the product \(e(v'/v)f(v'/v)\) divides \(n = [\mathbf{L}:\mathbf{K}]\) and that the quotient \(n / e(v'/v)f(v'/v)\) is a power of the characteristic exponent of the residue field \(k\) of \(v\) . (Reduce it to the case where L is a Galois extension of K; use Exercises 6 (e) and 7 and Chapter V, §2, Theorem 2 and Proposition 5.)
\item
  Suppose further that \(n\) is not divisible by the characteristic exponent of the residue field \(k\) of \(v\) and that \(f(v' / v) = 1\) . Show that, for every integer \(m\) equal to the order of an element \(\gamma\) of the group \(\Gamma_{v'} / \Gamma_v\) there exists \(x \in \mathbf{L}\) such that the class of \(v'(x) \mod \Gamma_v\) is equal to \(\gamma\) and \(x^m \in \mathbf{K}\) .
\end{enumerate}

¶ 10. Let w be the 2-adic valuation on the 2-adic field \(Q_{2}\) ; on the field of rational functions \(\mathbf{Q}_2(\mathbf{X})\) consider the discrete valuation \(v\) extending \(w\) and such that

\[
v (a _ {0} + a _ {1} \mathrm{X} + \dots + a _ {n} \mathrm{X} ^ {n}) = \inf _ {i} (w (a _ {i}))
\]

(§ 10, no. 1, Lemma 1); let K be the completion of \(\mathbf{Q}_2(\mathbf{X})\) with respect to this valuation and let its valuation also be denoted by \(v\) . Let L be the splitting field of the polynomial \((\mathrm{Y}^2 - \mathrm{X})^2 - 2\) in K{[}Y{]} and let \(v'\) be the unique valuation on L extending \(v\) . Show that {[}L: K{]} = 8, \(e(v'/v) = 4\) , \(f(v'/v) = 2\) ; if \(k\) (resp. \(k'\) ) is the residue field of \(v\) (resp. \(v'\) ), \(k'\) is a radical extension of \(k\) of degree 2; show that there exists no sub-extension E of L such that {[}E: K{]} = 2 for which the residue field of the restriction of \(v'\) to E is equal to \(k'\) . (Of necessity \(\mathrm{E} = \mathrm{K}(\sqrt{\alpha})\) , where \(\alpha \in \mathrm{K}\) would not be a square in K but such that \(\alpha \equiv \mathrm{X}\) (mod. 2); express \(\sqrt{\alpha}\) with the aid of a suitable basis of L over \(\mathrm{K}(\sqrt{2})\) and observe that there exists in \(\mathrm{K}(\sqrt{2})\) no element whose square is congruent to X mod. 2.)

¶ 11. Let K be a field, L a finite Galois extension of K and G its Galois group. Let v be a valuation on L with residue field k and order group Γ; suppose that v is invariant under G and that the restriction v\textbar K has the same residue field k as v, so that \(G = G^{Z} = G^{T}\) (notation of Exercise 8).

\begin{enumerate}
\def\labelenumi{(\alph{enumi})}
\item
  Let \(x \in \mathbf{L}^*\) , \(\sigma \in \mathcal{G}\) . Show that \(x^{-1}\sigma(x)\) is a unit for \(v\) , that its image \(\varepsilon_{\sigma}(x)\) in \(k^*\) depends only on the valuation \(v(x)\) and the class of \(\sigma\) modulo the commutator group \(\mathcal{G}'\) of \(\mathcal{G}\) and that there is thus defined a Z-bilinear mapping (also denoted by \(\varepsilon\) ) \((\mathcal{G}/\mathcal{G}') \times \Gamma \to k^*\) , which is equal to 1 on \((\mathcal{G}/\mathcal{G}') \times v(\mathbf{K}^*)\) .
\item
  For all \(\sigma \in \mathcal{G}\) , let \(\phi(\sigma)\) be the homomorphism \(\Gamma / v(\mathbf{K}^*) \to k^*\) mapping \(\varepsilon_{\sigma}(x)\) to \(v(x)\) for all \(x \in \mathbf{L}^*\) ; thus there is defined a homomorphism
\end{enumerate}

\[
\phi \colon \mathcal {G} \to \operatorname{Hom} _ {\mathbf {Z}} (\Gamma / v (\mathrm{K} ^ {*}), k ^ {*}).
\]

Show that, if \(p\) is the characteristic exponent of \(k\) , the kernel \(\mathcal{N}\) of \(\phi\) has order a power of \(p\) . (In the contrary case, there would exist a \(\sigma \neq e\) in \(\mathcal{N}\) and a prime number \(q \neq p\) such that \(\sigma^q = e\) ; for some integer \(x\) in \(L - K\) , write \(\sigma(x) = x + y\) , where \(v(y) > v(x)\) , and calculate \(\sigma^q(x)\) to obtain a contradiction.) Deduce that \(\mathcal{G}\) is a solvable group.

\begin{enumerate}
\def\labelenumi{(\alph{enumi})}
\setcounter{enumi}{2}
\tightlist
\item
  Suppose that \(n = [\mathbf{L} : \mathbf{K}]\) is prime to \(p\) ; show that \(\phi\) is bijective, that \(e(\mathbf{L} / \mathbf{K}) = n\) and that, if \(\nu\) is the \(lcm\) of the orders of the elements of \(\mathcal{G}\) , \(k^*\) contains the \(\nu\) -th roots of unity. (Use (b) and Lemma 2 of no. 1.)
\end{enumerate}

¶ 12. Let K be a field, v a valuation on K, A the ring of the valuation v and Γ its order group.

\begin{enumerate}
\def\labelenumi{(\alph{enumi})}
\tightlist
\item
  Let \(f(\mathbf{X}) = a_{0}\mathbf{X}^{n} + a_{1}\mathbf{X}^{n-1} + \cdots + a_{n}\) be a polynomial in A{[}X{]} such that \(v(a_{0}) = 0\) , all of whose roots belong to K; let \(x_{i} (1 \leqslant i \leqslant r)\) be these distinct roots and let \(k_{i}\) be the order of multiplicity of \(x_{i} (1 \leqslant i \leqslant r)\) . Let
\end{enumerate}

\[
g (\mathbf {X}) = b _ {0} \mathbf {X} ^ {n} + b _ {1} \mathbf {X} ^ {n - 1} + \dots + b _ {n} = b _ {0} (\mathbf {X} - y _ {1}) \dots (\mathbf {X} - y _ {n})
\]

another polynomial in A{[}X{]} such that \(v(b_{0}) = 0\) and the \(y_{h}\) belong to K. Show that there exists \(\lambda_{0} \in \Gamma\) such that \(v(x_{i} - x_{j}) < \lambda_{0}\) for every ordered pair of distinct indices i, j and, for all \(\lambda \geqslant \lambda_{0}\) , there exists \(\mu \geqslant \lambda\) such that the relations \(v(a_{i} - b_{i}) \geqslant \mu\) for all i imply that, for \(1 \leqslant i \leqslant r\) , there are exactly \(k_{i}\) indices h such that \(v(x_{i} - y_{i}) \geqslant \lambda\) . (First evaluate \(v(f(y_{h}))\) in two ways to show that necessarily \(v(x_{i} - y_{h}) \geqslant \lambda\) for at least one i; then evaluate

\[
v \left(\frac {f ^ {\prime} (z)}{f (z)} - \frac {g ^ {\prime} (z)}{g (z)}\right)
\]

at suitable points \(z \in \mathbf{K}\) .) If further \(b_{i} = a_{i}\) except for \(i = n - 1\) , show that, so long as \(\lambda_0\) is sufficiently large, the \(y_{h}\) are all distinct (evaluate \(f'(y_h)/f(y_h)\) assuming that \(y_{h}\) is not a simple root of \(g\) ).

\begin{enumerate}
\def\labelenumi{(\alph{enumi})}
\setcounter{enumi}{1}
\tightlist
\item
  Suppose henceforth that \(\mathbf{K}\) is the algebraic closure of a subfield \(\mathbf{K}_0\) such that \(\mathbf{A}_0 = \mathbf{A} \cap \mathbf{K}_0\) is Henselian. Suppose that the polynomial \(f\) belongs to \(\mathbf{A}_0[\mathbf{X}]\) and is irreducible and separable over \(\mathbf{K}_0\) . Let \(z \in \mathbf{K}\) be such that
\end{enumerate}

\[
v (z - v _ {i}) > v (z - x _ {j})
\]

for all \(j \neq i\) ; then show that \(\mathrm{K}_0(x_i) \subset \mathrm{K}_0(z)\) . (Show that \(z - x_i\) is equal to all its conjugates with respect to \(\mathrm{K}_0(x_i)\) .)

\begin{enumerate}
\def\labelenumi{(\alph{enumi})}
\setcounter{enumi}{2}
\tightlist
\item
  Suppose that \(f \in \mathbf{A}_0[\mathbf{X}]\) is separable and let \(f = a_0 \prod_{i=1}^{\infty} f_i\) be the decomposition of \(f\) into monic irreducible polynomials in \(\mathbf{K}_0[\mathbf{X}]\) (which belong in fact to \(\mathbf{A}_0[\mathbf{X}]\) , cf.~Chapter V, § 1, no. 3, Proposition 11). Show (in the notation of (a)) that, if \(\mu\) is taken sufficiently large, the decomposition of \(g\) into monic
\end{enumerate}

irreducible factors in \(K_{0}[X]\) can be written as \(b_{0}\prod_{i=1}^{n}g_{i}\) , where, for all i, \(g_{i}\) is of the same degree as \(f_{i}\) (a decomposition said to be ``of the same type'' as that of f) and further the splitting fields of \(f_{i}\) and \(g_{i}\) are the same for all i. (Show first that g is separable; consider then a finite Galois extension of \(K_{0}\) containing the roots of f and g and consider the way in which the Galois group of this extension permutes these roots; finally use (b).)

\begin{enumerate}
\def\labelenumi{(\alph{enumi})}
\setcounter{enumi}{3}
\tightlist
\item
  Give an example where \(f\) is irreducible but not separable and where, for all \(\mu \geqslant \lambda_0\) , there exists a polynomial
\end{enumerate}

\[
g (\mathbf {X}) = b _ {0} \mathbf {X} ^ {n} + b _ {1} \mathbf {X} ^ {n - 1} + \dots + b _ {n} \in \mathrm{A} _ {0} [ \mathbf {X} ]
\]

where \(v(a_{i} - b_{i}) \geqslant \mu\) for all i, which is not irreducible (take \(K_{0}\) such that \(K \neq K_{0}\) and K is a radial extension of \(K_{0}\) ).

¶ 13. (a) Let K be a field and v a non-improper valuation on K. Let E be a set filtered by an ultrafilter U and \(\xi \mapsto x_{\xi}\) , \(\xi \mapsto y_{\xi}\) two mappings of E to K with the topology \(\mathcal{T}_{v}\) . Show that, if the mapping \(\xi \mapsto x_{\xi}y_{\xi}\) converges to 0 in K with respect to U, so does one of the mappings \(\xi \mapsto x_{\xi}\) , \(\xi \mapsto y_{\xi}\) . (Observe that, if \(\xi \mapsto x_{\xi}\) has no cluster point with respect to \(\mathfrak{U}\) , the mapping \(z \mapsto x_{\xi}^{-1}\) is bounded on a subset of \(\mathfrak{U}\) .)

\begin{enumerate}
\def\labelenumi{(\alph{enumi})}
\setcounter{enumi}{1}
\item
  Let \(f\) be a non-constant polynomial in \(\mathbf{K}[\mathbf{X}]\) . Show that, if \(\xi \mapsto f(x_{\xi})\) converges to 0 in \(\mathbf{K}\) with respect to \(\mathfrak{U}\) , \(\xi \mapsto x_{\xi}\) converges in \(\hat{\mathbf{K}}\) (decompose \(f\) into factors in an algebraic extension of \(\mathbf{K}\) and use (a)).
\item
  Suppose that K is algebraically closed in \(\hat{\mathbf{K}}\) . Then the mapping \(x \mapsto f(x)\) of K to itself is closed (use (b)). Deduce that, if \(g\) is a rational function in \(\mathbf{K}(\mathbf{X})\) and B is a bounded closed subset of K, then \(g(\mathbf{B} - \mathbf{P})\) (where P is the set of poles of \(g\) in B) is closed in K.
\item
  Suppose that the algebraic closure of \(\mathbf{K}\) is a radicial extension of \(\mathbf{K}\) . Then show that \(\hat{\mathbf{K}}\) is algebraically closed. (Apply (b) to an algebraic closure of \(\hat{\mathbf{K}}\) and a polynomial \(f \in \hat{\mathbf{K}}[\mathbf{X}]\) ; use also Exercise 12 (a).)
\end{enumerate}

¶ 14. (a) For a field K to be Henselian with a valuation v, it is necessary and sufficient that K be the greatest separable algebraic extension of K contained in \(\hat{K}\) and that \(\hat{K}\) be Henselian. (To see that the condition is necessary, use Exercise 12 (b); to show that it is sufficient, observe that, to verify condition (H) of Chapter III, § 4, Exercise 3, we may confine our attention to the case where P is separable over K, noting that, if Q is an irreducible factor of P in \(\hat{K}[X]\) such that \(Q^{p^{e}}\) belongs to K{[}X{]}, then Q is a g.c.d. of P and \(Q^{p^{e}}\) in \(\hat{K}[X]\) .) Consider the case where v is of height 1 (cf.~Exercise 6 (b)).

\begin{enumerate}
\def\labelenumi{(\alph{enumi})}
\setcounter{enumi}{1}
\item
  Deduce from (a) that there exists in the \(p\) -adic field \(\mathbf{Q}_p\) a countable subfield which is Henselian with the \(p\) -adic valuation.
\item
  Give an example of a field K which is Henselian with a discrete valuation which is not complete and such that \(\hat{K}\) is a finite radical extension of K (cf.~Chapter V, § 1, Exercise 20).
\end{enumerate}

¶ 15. (a) Let K be a field, \(v_1\) , \(v_2\) two independent valuations (§ 7, no. 2) on K and \(K_1\) , \(K_2\) the completions of K with respect to \(v_1\) and \(v_2\) respectively. Let \(L_1\) , \(L_2\) be two fields such that \(K \subset L_1 \subset K_1\) , \(K \subset L_2 \subset K_2\) , which are Henselian (with the extensions by continuity of \(v_1\) and \(v_2\) respectively). Let \(g_1 \in L_1[X]\) , \(g_2 \in L_2[X]\) be two separable polynomials with the same degree n.~Show that there exists a polynomial \(h \in K[X]\) which, in \(L_i[X]\) , has the same type of decomposition (Exercise 12 (c)) as \(g_i\) ( \(i = 1, 2\) ). (Reduce it to the case where \(g_1\) and \(g_2\) are in K{[}X{]} and are monic; consider the polynomial

\[
a ^ {n} g _ {1} (\mathbf {X} / a) + b ^ {n} g _ {2} (\mathbf {X} / a) - \mathbf {X} ^ {n},
\]

where \(v_{1}(a) \leqslant 0\) , \(v_{2}(a) > 0\) is arbitrarily large and \(v_{1}(b) > 0\) is arbitrarily large.)

\begin{enumerate}
\def\labelenumi{(\alph{enumi})}
\setcounter{enumi}{1}
\item
  Deduce from (a) that, if K is Henselian with \(v_{1}\) , the algebraic closure of \(L_{2}\) is radial over \(L_{2}\) and \(K_{2}\) is algebraically closed (take \(g_{2}\) to be irreducible and \(g_{1} \in K[X]\) to be the product of n distinct factors of prime degree).
\item
  Deduce from (b) that, if \(\mathbf{K}\) is Henselian with \(v_{1}\) and \(v_{2}\) , the algebraic closure of \(\mathbf{K}\) is radicial over \(\mathbf{K}\) .
\item
  Show that, if K is Henselian with a discrete valuation v, the algebraic closure of K cannot be radicial over K and therefore K cannot be Henselian with any valuation independent of v (use Algebra, Chapter V, § 11, Exercise 12).
\end{enumerate}

\begin{enumerate}
\def\labelenumi{\arabic{enumi}.}
\setcounter{enumi}{15}
\tightlist
\item
  Let K be a Henselian field with a valuation v of height 1. If L is an infinite separable algebraic extension of K, show that L cannot be complete with the valuation extending v (and also denoted by v). (Form a sequence \((x_{p})\) of elements of L such that the degree \(n_{p}\) of \(x_{p}\) with respect to K tends to \(+\infty\) and \(v(x_{p+1}-x_{p})\) is strictly greater than the valuations of the differences between \(x_{p}\) and its conjugates with respect to K; show that the sequence \((x_{p})\) is a Cauchy sequence which cannot converge in E, using Exercise 12 (b).)
\end{enumerate}

¶ 17. Let K be a Henselian field with a valuation v which is not algebraically closed, and L a subfield of K such that {[}K: L{]} is finite. Show that under these conditions L is Henselian with the restriction of v. (In the contrary case, there would exist two extensions \(v_{1}, v_{2}\) of \(v|L\) to an algebraic closure \(\Omega\) of K such that the restriction of \(v_{1}\) and \(v_{2}\) to a finite algebraic extension \(L'\) of L would not be equivalent. Deduce that there would exist on a suitable finite quasi-Galois extension \(K'\) of L containing K, two valuations \(v_{1}', v_{2}'\) which are not equivalent and for which \(K'\) would be Henselian; using Exercise 15 (c), show that, if \(E = L^{p^{-\infty}}\) (p being the characteristic exponent of \(\Omega\) ), \(E \neq \Omega\) but \(\Omega\) would be a finite extension of E; complete the argument with the aid of Algebra, Chapter VI, § 2, Exercise 31.)

¶ 18. Let K be a Henselian field with a valuation v, k the residue field of v and suppose that every finite algebraic extension of k is cyclic over k (which is for example the case when k is finite). Let A be the ring of v and \(f \in A[X]\) a monic polynomial such that, if \(f(\mathbf{X}) = (\mathbf{X} - \alpha_{1}) \ldots (\mathbf{X} - \alpha_{n})\) , where the \(\alpha_{i}\) belong to an algebraic closure of K, the element \(D = \prod_{i < j} (\alpha_{i} - \alpha_{j})^{2}\) of A (``discriminant'' of f) is such that \(v(D) = 0\) . Let \(\bar{f}_{j} (1 \leqslant j \leqslant s)\) be the irreducible factors of the canonical image \(\bar{f}\) of f in k{[}X{]} and let \(r_{j}\) be the degree of \(\bar{f}_{j}\) . Show that the Galois over K of the splitting field of f, considered as a group of permutations of the \(\alpha_{i}\) is generated by a permutation \(\sigma\) which decomposes into cycles of respective lengths \(r_{1}, r_{2}, \ldots, r_{s}\) (observe that there exists only a single extension of k of given degree, up to a k-isomorphism). Deduce that, for D to be a square in A, it is necessary and sufficient that n - s be even (``Stickelberger's Theorem''; examine under what condition D is invariant under \(\sigma\) ).

\begin{enumerate}
\def\labelenumi{\arabic{enumi}.}
\setcounter{enumi}{18}
\tightlist
\item
  Let K be a Henselian field with a valuation v. Let E be a finite-dimensional vector space over K and Q a non-degenerate quadratic form on E such that the relation \(Q(x) = 0\) implies x = 0. Let
\end{enumerate}

\[
\Phi (x, y) = \mathrm{Q} (x + y) - \mathrm{Q} (x) - \mathrm{Q} (y)
\]

be the associated bilinear form. Show that \(2v(\Phi(x,y)) \geqslant v(Q(x)) + v(Q(y))\) ; deduce that

\[
v (\mathrm{Q} (x + y)) \geqslant \inf (v (\mathrm{Q} (x)), v (\mathrm{Q} (y))).
\]

¶ 20. Let K be a field of characteristic ≠2 which is Henselian with a valuation v; let ξ ↦ ξ̄ be an involutive automorphism of K such that v(ξ̄) = v(ξ) and let Φ be a non-degenerate Hermitian form on a finite-dimensional vector space E over K.

\begin{enumerate}
\def\labelenumi{(\alph{enumi})}
\item
  Let \(U = (\alpha_{ij})\) be the matrix of \(\Phi\) with respect to a basis \((e_i)\) of E; show that there exists an element \(\lambda\) of the order group of \(v\) such that, if \(\Phi'\) is another Hermitian form on E whose matrix \(U' = (\alpha_{ij}')\) with respect to \((e_i)\) satisfies the conditions \(v(\alpha_{ij} - \alpha_{ij}') > \lambda\) for every ordered pair \((i,j)\) , then \(\Phi\) and \(\Phi'\) are equivalent. (Argue as in Algebra, Chapter IX, § 6, Exercise 6, using Exercise 14 (a) above.)
\item
  Suppose that \(\xi \mapsto \overline{\xi}\) is the identity (hence \(\Phi\) is a symmetric form) and that the valuation \(v\) is discrete, normed and such that \(v(2) = 0\) ; let \(\pi\) be a uniformizer for \(v\) . Show that \(\Phi\) is characterized, up to equivalence, by its index \(\nu\) and by two symmetric bilinear forms \(\Psi_1, \Psi_2\) of index 0 on vector spaces \(k^r\) and \(k^s\) respectively, where \(k\) is the residue field of \(v\) and \(r + s = n - 2\) . (With the aid of a Witt decomposition, reduce it to the case where \(\Phi\) is of index 0; with the aid of Exercise 19, show that, for all \(i \geqslant 0\) , the set \(\mathbf{M}_i\) of \(x \in \mathbf{E}\) such that \(v(\Phi(x, x)) \geqslant i\) is a module over the ring A of the valuation \(v\) . Show that, if \(x \in \mathbf{M}_i\) , \(y \in \mathbf{M}_{i+1}\) , then \(v(\Phi(x, y)) \geqslant i + 1\) and, by taking quotients, symmetric bilinear forms can therefore be derived from \(\Phi\) on the vector \(k\) -spaces \(\mathbf{M}_0 / \mathbf{M}_1\) and \(\mathbf{M}_1 / \mathbf{M}_2\) . Use Exercise 6 (c) of §3 and the fact that the equation \(\xi^2 = \alpha\) has a solution in A for \(\alpha \equiv 1 (\text{mod. } \pi)\) .) Consider the case where \(k\) is finite (cf.~Algebra, Chapter IX, §6, Exercise 4).
\end{enumerate}

\begin{enumerate}
\def\labelenumi{\arabic{enumi}.}
\setcounter{enumi}{20}
\tightlist
\item
  Let K be a field, v a discrete valuation on K, A the ring of v, \(\pi\) a uniformizer of v and k the residue field of v.
\end{enumerate}

\begin{enumerate}
\def\labelenumi{(\alph{enumi})}
\item
  Let P, R be two polynomials in A{[}X{]}, P being monic, and suppose that: (1) \(\deg(R) < h \cdot \deg(P)\) , where \(h \geqslant 1\) is an integer; (2) the canonical image \(\overline{P}\) of P in k{[}X{]} is irreducible. Show then that, if the polynomial \(Q = P^{h} + \pi R\) is reducible in \(\widehat{K}[X]\) , necessarily h \textgreater{} 1 and \(\overline{P}\) divides the canonical image \(\overline{R}\) of R in k{[}X{]}. Deduce that, for a given monic irreducible polynomial \(r(X) \in k[X]\) and a given integer \(h \geqslant 1\) , there exists a separable irreducible polynomial \(Q \in A[X]\) such that its canonical image \(\overline{Q}\) is equal to \(r^{h}\) .
\item
  Let \(k' = k(\alpha)\) be an algebraic extension of \(k\) of degree \(m\) and \(h\) an integer \(\geqslant 1\) . Show that there exists an algebraic extension \(L\) of \(K\) of degree \(hm\) such that there is only a single valuation \(v'\) (up to equivalence) on \(L\) extending \(v\) , that \(e(v'/v) = h\) , \(f(v'/v) = m\) and that the residue field of \(v'\) is isomorphic to \(k'\) (use (a)).
\end{enumerate}

\begin{enumerate}
\def\labelenumi{\arabic{enumi}.}
\setcounter{enumi}{21}
\tightlist
\item
  \begin{enumerate}
  \def\labelenumii{(\alph{enumii})}
  \tightlist
  \item
    If there exists a discrete valuation on a field \(\mathbf{K}\) , show that the algebraic closure of \(\mathbf{K}\) is infinite over \(\mathbf{K}\) .
  \end{enumerate}
\end{enumerate}

\begin{enumerate}
\def\labelenumi{(\alph{enumi})}
\setcounter{enumi}{1}
\tightlist
\item
  Let K be a finitely generated extension of a field \(K_{0}\) . Show that, if K is not algebraic over \(K_{0}\) , there exists a discrete valuation v on K such that \(v(x) = 0\) on \(K_{0}\) .
\end{enumerate}

